\chapter{Исходные коды измерительных скриптов}
\label{sec:Appendix1}

В оригинале работы (приложение A) напечатаны исходные коды семи ключевых скриптов: \texttt{setup\_network.sh}, \texttt{sniff\_stats.py}, \texttt{pkt\_gen.py}, \texttt{perf\_test.py}, \texttt{ra\_audit.py}, \texttt{neigh\_scan.py} и \texttt{arp\_scan.py}. В этом переводе коды не повторяются: они одинаковы в обеих версиях и лежат в папке \texttt{SourceCodes/} проекта.

\chapter{Обзор программного обеспечения прибора}
\label{sec:Appendix2}

Программное обеспечение хранится в каталоге \texttt{\textasciitilde/diploma\_project/} на карте памяти прибора; системные файлы лежат в обычных местах системы. Полные исходные коды всех файлов входят в электронное приложение к работе. Обзор приведён в таблице \ref{tab:soubory}; скрипты, напечатанные в приложении A оригинала, отмечены звёздочкой.

\begin{table}[htbp]
    \caption{Файлы программного обеспечения прибора}
    \label{tab:soubory}
    \centering
    \small
    \begin{tabular}{lp{9.5cm}}
        \hline
        Файл & Назначение \\
        \hline
        \multicolumn{2}{l}{\textit{Измерительные скрипты, запускаются в пространстве имён порта}} \\
        \texttt{setup\_network.sh}* & подготовка портов, пространств имён и сервера dnsmasq \\
        \texttt{sniff\_stats.py}* & анализатор трафика, статистика, сохранение в pcap и CSV \\
        \texttt{pktdecode.py} & разбор заголовков кадра без Scapy, для статистики и для разбора по слоям \\
        \texttt{pkt\_gen.py}* & генератор пакетов шести типов \\
        \texttt{perf\_test.py}* & измерения iperf3 и развёртка по размерам кадра \\
        \texttt{perf\_server.py} & сервер iperf3 на выбранном порту \\
        \texttt{path\_test.py} & тест доступности и задержки утилитой ping \\
        \texttt{targets.py} & общий поиск цели и объяснение, почему цели нет \\
        \texttt{ra\_audit.py}* & аудит сообщений Router Advertisement \\
        \texttt{neigh\_scan.py}* & поиск соседей IPv6 \\
        \texttt{arp\_scan.py}* & сканирование ARP \\
        \multicolumn{2}{l}{\textit{Управляющее приложение}} \\
        \texttt{gui\_app.py} & главный экран, запуск задач, экран приветствия и блокировки \\
        \texttt{ui\_kit.py} & общий вид экранов, строка состояния, кнопки \\
        \texttt{screens.py} & экраны PATH, SCAN, IPv6, DHCP / RA и SYSTEM \\
        \texttt{traffic.py} & экран TRAFFIC с разбором кадра и фильтрами \\
        \texttt{generator.py} & экран GENERATOR \\
        \texttt{speed.py} & экран SPEED с развёрткой и режимом сервера \\
        \texttt{results.py} & сохранение результатов в файлы JSON \\
        \texttt{wifi\_dialog.py} & подключение к сети Wi-Fi и клавиатура на экране \\
        \multicolumn{2}{l}{\textit{Системные службы}} \\
        \texttt{port\_status.py} & слежение за состоянием портов и запись в файл состояния \\
        \texttt{ups\_monitor.py} & измерение состояния аккумулятора и штатное выключение при разряде \\
        \texttt{analyzer-port@.service} & подготовка порта при подключении адаптера \\
        \texttt{10-analyzer-*.link} & постоянные имена портов по физическому адресу \\
        \texttt{\textasciitilde/.xinitrc} & запуск управляющего приложения вместо рабочего стола \\
        \hline
    \end{tabular}
\end{table}
